% =============================================================================
% AGENTE  : Statistician (outputs) + Geneticista (interpretação)
% MÓDULO  : Resultados — Português Brasileiro (results_pt.tex)
% SINCRON.: Deve ser consistente com results_en.tex
% STATUS  : RASCUNHO
% =============================================================================

\section{Resultados}
\label{sec:results}

\subsection{Visão Geral do Painel de Marcadores Simulados}

Um total de 60 marcadores moleculares foi avaliado, compreendendo 36 \snp{} e
24 \ssr{}. As estatísticas descritivas para \pic{}, \he{} e \ho{} são
apresentadas na \cref{tab:marker_summary} (Material Suplementar).

\subsection{Conteúdo de Informação sobre Polimorfismo}

Os marcadores \snp{} apresentaram média de \pic{} de aproximadamente
0,25--0,32, refletindo a restrição inerente da variação bialélica. Em
contraste, os marcadores \ssr{}, modelados com distribuições de frequência
alélica simétricas, exibiram valores médios de \pic{} substancialmente mais
elevados (0,35--0,42), consonantes com seu reconhecido papel como sistemas
de marcadores altamente polimórficos \autocite{Powell1996}.

A distribuição dos valores de \pic{} está ilustrada na
\cref{fig:pic_boxplot}. Os marcadores \ssr{} apresentaram amplitude
interquartílica mais ampla, indicando maior variabilidade de informatividade
entre os locos.

\subsection{Heterozigose}

A heterozigose esperada (\he{}) seguiu o mesmo padrão do \pic{}, uma vez que
ambas as estatísticas são funções diretas das frequências alélicas. O valor
médio de \he{} para os marcadores \ssr{} superou o dos marcadores \snp{},
refletindo as frequências alélicas mais equilibradas dos primeiros. Os
valores de heterozigose observada (\ho{}) distribuíram-se próximos a
\he{} (desvio médio $\leq 0,04$), sugerindo ausência de desvios
relevantes do equilíbrio de Hardy--Weinberg na população simulada.
